\chapter{مبانی نظری و مدل‌سازی سیستم}
\section{مقدمه}

هدف این فصل، ارائه‌ی چارچوب نظری مستحکمی‌است که درک صحیح راهکار پیشنهادی \lr{RS-DRL} را ممکن می‌سازد. سیستم‌های نرم‌افزاری مدرن همواره با عدم قطعیت‌های مداوم مواجه هستند و باید در زمان اجرا به‌طور پیوسته ساختار و رفتار خود را با شرایط در حال تغییر وفق دهند. این فصل در ابتدا اصول اولیه سیستم‌های خودوفقی و مؤلفه‌های آن را تشریح می‌کند، سپس منابع عدم قطعیت و روش‌های مدل‌سازی آن را معرفی می‌نماید. در ادامه، با معرفی سیستم \lr{DeltaIoT} به‌عنوان یک مطالعه موردی عینی، مبانی یادگیری تقویتی و الگوریتم \lr{DQN} را در قالب این سیستم فرموله کرده و در نهایت مسئله اصلی رساله را به صورت صوری بیان می‌کنیم.

\section{مبانی سیستم‌های خودوفقی}
\subsection{مفاهیم اصلی و الگوهای معماری}

سیستم‌های خودوفقی (\lr{Self-Adaptive Systems}) یک رویکرد برجسته و مؤثر برای مقابله با عدم قطعیت به شمار می‌روند. ایده کلیدی در این سیستم‌ها، تجهیز آن‌ها به یک یا چند حلقه بازخوردی کنترلی است که به صورت مداوم مؤلفه‌های داخلی سیستم و محیط خارجی آن را پایش کرده و در صورت نیاز سیستم را با تغییرات منطبق می‌نمایند \cite{RS_DRL_Article}. 

مهمترین و شناخته‌شده‌ترین الگوی معماری برای پیاده‌سازی این سیستم‌ها، حلقه خودمختار \lr{MAPE-K} است. این حلقه در کنار سیستم هدف و منطق کسب‌وکار، جریان کنترلی زیر را اعمال می‌کند:
\begin{itemize}
  \item \textbf{پایش (\lr{Monitor}):} سنسورها وضعیت محیط فیزیکی، زیرساخت و خود سیستم را به صورت زمان‌حقیقی استخراج می‌کنند.
  \item \textbf{تحلیل (\lr{Analyze}):} داده‌های پایش‌شده را پردازش کرده و مشخص می‌کند که آیا انحرافی از اهداف کیفی و کمی سیستم رخ داده است و نیاز به اعمال تغییر وجود دارد یا خیر.
  \item \textbf{برنامه‌ریزی (\lr{Plan}):} در صورت نیاز تطبیق، راهکارها و استراتژی‌های موجود (گزینه‌های وفقی) مقایسه شده و گزینه بهینه برای شرایط فعلی برگزیده می‌شود.
  \item \textbf{اجرا (\lr{Execute}):} استراتژی برنامه‌ریزی‌شده در قالب یک سری از دستورات اجرایی از طریق عملگرها (Actuators) بر روی سیستم اعمال می‌گردد.
  \item \textbf{دانش (\lr{Knowledge}):} به‌عنوان حافظه مرکزی عمل کرده و اطلاعات، مدل‌ها و داده‌های بین چهار مرحله قبلی را هماهنگ و نگهداری می‌کند. راهکار پیشنهادی ما (\lr{RS-DRL}) مستقیماً به‌عنوان هسته اصلی بخش برنامه‌ریزی در این الگو فعالیت می‌کند.
\end{itemize}

\subsection{چالش‌های عدم قطعیت در زمان اجرا}

عدم قطعیت در سیستم‌های خودوفقی، به هرگونه انحراف از دانش قطعی اطلاق می‌شود که می‌تواند باعث کاهش اطمینان نسبت به تصمیمات زمان اجرا گردد \cite{RS_DRL_Article}. 

به عنوان نمونه‌های عینی، می‌توان به موارد گوناگونی در دامنه‌های مختلف اشاره کرد: در برنامه‌های مبتنی بر اینترنت اشیاء (\lr{IoT}) که در محیط‌های شهری پراکنده هستند، تغییرات مداوم و غیرقابل پیش‌بینی آب و هوا منجر به افت کیفیت سیگنال و اختلال در ارسال اطلاعات سنسورها می‌شود. در محیط‌های ابری، نوسان شدید در تقاضای منابع از سوی کاربران، و در حوزه‌های امنیتی، ظهور بی‌وقفه و تصادفی حملات سایبری از مصادیق برجسته عدم قطعیت هستند.

اهمیت و چالش‌برانگیز بودن عدم قطعیت از آنجا نشأت می‌گیرد که مفروضات، پیکربندی‌ها و تصمیمات در نظر گرفته شده در فاز طراحی معماری نرم‌افزار را بی‌اعتبار می‌سازد. در نتیجه، سیستم باید قادر باشد با یادگیری آنلاین و پویا در زمان اجرا، اهداف اصلی خود را محقق نماید؛ وگرنه با نقض کیفیت خدمات صدمات سنگین احتمالی را تجربه خواهد نمود.

\section{مدل‌سازی عدم قطعیت در سیستم‌های خودوفقی}
\subsection{منابع عدم قطعیت}

منابع عدم قطعیت به‌طور کلی در سه دسته وسیع طبقه‌بندی می‌شوند: نخست، نامشخص بودن موجودیت‌های محیطی نظیر پارازیت در حسگرها، نویزها و تغییرات جوی که از کنترل ساختار درونی سیستم خارج هستند. دوم، ماهیت غیرقابل اتکای منابع خود سیستم مثل نقص در قطعات زیربنایی، کاهش سریع سطح باتری یا مشکلات نرم‌افزاری تعاملات میان سرویس‌ها؛ و در نهایت عدم قطعیت در اهداف متغیر و ترجیحات نامشخص ذینفعان در محیط‌های تجاری.

\subsection{رویکردهای مدل‌سازی در زمان اجرا}

جهت درک و محاسبه رفتار سیستم در هنگام بروز مقادیر ناشناخته، استفاده از مدل‌سازی‌های زمان اجرا (\lr{Runtime Models}) ضروری به نظر می‌رسد. در شرایطی که فضای جستجو بسیار بزرگ و مخاطره تصمیم‌گیری به دلیل شرایط مبهم بالاست، رویکرد پیشنهادی از مکانیزم‌های مدل‌سازی ریاضی صوری بهره‌گیری و پشتیبانی می‌کند.

در این راستا ابزارهایی مانند موتور \lr{UPPAAL-SMC} و چارچوب \lr{ActivFORMS} به کار گرفته می‌شوند. این زیرساخت‌ها با استفاده از اتوماتای زمان‌دار (\lr{Timed Automata}) مدلی رفتاری از سیستم هدف و محیط آن را توسعه داده و از طریق روش‌های کنترل مدل آماری (\lr{Statistical Model Checking}) می‌توانند به سرعت گزینه‌های وفقی پیشنهاد شده توسط سیستم هوشمند را پیش از اجرا در سیستم فیزیکی ارزیابی، تایید یا رد کنند. این قابلیت نقشی کلیدی در امنیت راهکار پیشنهادی (در بخش تاییدگر گزینه وفقی) این رساله ایفا می‌کند.

\subsection{مطالعه موردی انگیزشی: سیستم \lr{DeltaIoT}}

برای عینی‌تر ساختن این مفاهیم انتزاعی در یک محیط عملیاتی، در ادامه سیستم \lr{DeltaIoT} را معرفی می‌کنیم؛ یک شبکه حسگر بی‌سیم که به‌عنوان یک معیار استاندارد، چالش‌های اصلی فضاهای وفقی بزرگ و عدم قطعیت‌های محیطی را به خوبی به تصویر می‌کشد.

\subsubsection{مرور کلی سیستم و منابع عدم قطعیت}

سیستم \lr{DeltaIoT} یک شبکه حسگر بی‌سیم (\lr{WSN}) با تعاملات پیچیده است که به صورت گسترده به‌عنوان یک محک استاندارد (\lr{Benchmark}) برای ارزیابی سیستم‌های خودوفقی به کار می‌رود \cite{RS_DRL_Article}. این شبکه بسته به نسخه مورد استفاده، مابین ۱۵ تا ۳۷ گره (سنسور) دارد. دو منبع اصلی عدم قطعیت در این سیستم در زمان اجرا بروز می‌یابند:
\begin{enumerate}
  \item \textbf{نویز و تداخل شبکه (\lr{Network Interference/Noise}):} تغییرات محیطی و مداخلات پیرامونی که موجب نوسانات پیش‌بینی‌نشده در کیفیت ارتباطات از \lr{-40dB} تا \lr{+15dB} می‌گردد.
  \item \textbf{تغییرپذیری بار پیام‌ها (\lr{Message Load Variability}):} نوسان در تعداد پیام‌های تولید شده توسط هر سنسور (بین ۰ تا ۱۰ پیام) با الگوهای انتقال متفاوت در طول زمان.
\end{enumerate}

همانطور که در مطالعات پیشین (و روند نشان داده شده در شکل ۲ مقاله مرجع) نمایان است، این عدم قطعیت‌ها باعث شکل‌گیری افت شدید عملکرد سیستم در طی چرخه‌های وفقی گشته و اهمیت حیاتی مکانیسم‌های تطابق پویا را برجسته می‌کنند.

\subsubsection{فضای وفقی و بازنمایی پیکربندی}

فضای وفقی (\lr{Adaptation Space}) دربرگیرنده مجموعه تمامی‌گزینه‌ها و پیکربندی‌های ممکنی است که سیستم می‌تواند در زمان اجرا اتخاذ نماید. در \lr{DeltaIoT}، سه نسخه متفاوت با مشخصاتی به شرح زیر طراحی گردیده است \cite{RS_DRL_Article}:
\begin{itemize}
  \item \textbf{\lr{DeltaIoTv1}:} شامل ۱۵ سنسور و دارای ۲۱۶ گزینه وفقی متمایز.
  \item \textbf{\lr{DeltaIoTv1.1}:} شامل ۱۶ سنسور و دربرگیرنده ۱۲۹۶ گزینه وفقی است.
  \item \textbf{\lr{DeltaIoTv2}:} نسخه وسیع‌تر این شبکه متشکل از ۳۷ سنسور که دارای فضای عظیم متشکل از ۴۰۹۶ تصمیم ممکن می‌باشد.
\end{itemize}

هریک از این گزینه‌های وفقی در حقیقت ترکیبی از پارامترهای مختلف هستند؛ نظیر سطوح توان انتقال انرژی (\lr{Transmission Power}) برای سنسورهای مشخص و توزیع نسبی مسیرهای شبکه‌ای (\lr{Link Distribution}). انتخاب یک ترکیب از این پارامترها، یک پیکربندی کاملاً جدید و واحد را شکل می‌دهد.

\section{یادگیری تقویتی برای وفقی‌سازی}

\subsection{اصول یادگیری تقویتی (\lr{MDP}، یادگیری \lr{Q})}

یادگیری تقویتی از طریق تعامل مستمر یک عامل (\lr{Agent}) با محیط اطرافش بر پایه آزمون و خطا بنا شده است و این محیط به لحاظ ریاضی از طریق فرآیند تصمیم‌گیری مارکوف (\lr{MDP}) فرموله می‌شود. فرآیند \lr{MDP} خود شامل فضای حالت (\lr{State})، فضای عمل (\lr{Action}) و توابع پاداش (\lr{Reward}) است.

فرآیند تصمیم‌گیری مارکوف، برای هر حالت خاص، عامل عمل مشخصی را اتخاذ می‌کند و مطابق با معادله مشهور بلمن (\lr{Bellman Equation}) محیط ارزش عمل اتخاذ شده را در قالب یک پاداش ارزیابی می‌نماید و به وضعیت بعدی منتقل می‌شود. این مفهوم در سیستم‌های خودوفقی مانند محیط \lr{DeltaIoT} به این صورت اعمال می‌گردد که وضعیت سیستم با معیارهایی نظیر انرژی ذخیره شده سیستم و تاخیر اندازه‌گیری شده فضای حالت را تشکیل می‌دهند و اقدام عامل، همان گزینه وفقی مورد انتخاب خواهد بود و محاسبه پاداش نیز طبق معادله بهینه‌سازی سیستم حاصل می‌گردد. نقش الگوریتم مشهور یادگیری Q ایجاد جدولی جامع برای تخصیص و پیدا کردن ارزش‌هاست تا مدل به سیاست بهینه همگرا شود.

\subsection{یادگیری تقویتی عمیق (\lr{DQN}) برای فضاهای بزرگ}

اما زمانی که با سیستم‌های واقعی که دارای فضای عمل گسسته ولی بسیار پهناور (برای مثال ۴۰۹۶ گزینه وفقی در شبکه‌های اینترنت اشیاء) مواجه هستیم، استفاده از متد کلاسیک یادگیری \lr{Q} با چالش نفرین ابعاد (\lr{Curse of Dimensionality}) روبه‌رو می‌شود؛ زیرا ذخیره و بروزرسانی جدولی با میلیون‌ها حالت و عمل از لحاظ محاسباتی ناکارآمد و در بسیاری موارد غیرممکن است.

شبکه‌های \lr{Q} عمیق (\lr{DQN}) راهکاری هوشمندانه برای این محدودیت ارائه می‌دهند؛ آن‌ها به جای جداول، تابع ارزش کیفیت \lr{Q} را با استفاده از شبکه‌های عصبی عمیق تقریب می‌زنند. برای غلبه بر بی‌ثباتی آموزش شبکه‌های عصبی در \lr{RL}، مدل \lr{DQN} از دو نوآوری برجسته بهره می‌برد: نخست استفاده از مکانیسم حافظه بازپخش تجربه (\lr{Experience Replay}) که نمونه‌های آموزشی تصادفی‌سازی شده و باعث شکستن همبستگی زمانی متوالی می‌گردد؛ و دوم تعبیه هدف ثابت به واسطه شبکه هدف ثابت (\lr{Target Network}) که بهبود تثبیت توابع هزینه را در هنگام محاسبه پیش‌بینی‌ها به دنبال دارد.

\subsection{مفهوم یادگیری مداوم در حلقه‌های وفقی}

اصلی‌ترین کاربرد این رساله، ادغام مفهوم آموزش پیوسته و مداوم توسط عامل \lr{DQN} در چارچوب معماری خودوفقی و به‌طور خاص درون حلقه \lr{MAPE-K} (همانطور که در فصل‌های آتی به تفصیل نشان داده خواهد شد) است. پس از جاگذاری در مؤلفه برنامه‌ریزی (\lr{Planner})، عامل یادگیرنده وضعیتی را به‌عنوان ورودی فضای حالت از ماژول پایش دریافت نموده، استراتژی اقدام را برنامه‌ریزی کرده و بر اساس میزان موفقیت انتخابش از سوی مؤلفه تحلیل پاداش می‌گیرد. این ساختار تعاملی، سیستم خودوفقی را از پیکربندی ایستا به ساختار یادگیری پویا تبدیل می‌سازد.

\subsection{مدل صوری برای تصمیم‌گیری تحت عدم قطعیت (ویژه \lr{DeltaIoT})}

با تبیین مبانی نظری فرآیند تصمیم‌گیری مارکوف و شبکه‌های \lr{DQN}، اکنون این چارچوب را برای سیستم \lr{DeltaIoT} پیاده‌سازی می‌کنیم. در ادامه، فضای حالت، فضای عمل و توابع پاداش اختصاصی که در طول این رساله مورد استفاده قرار می‌گیرند، تعریف خواهند شد.

برای امکان‌پذیر کردن یادگیری تقویتی، نیاز به فرمول‌بندی مسأله تطابق شبکه در قالب یک فرآیند تصمیم‌گیری مارکوف (\lr{MDP}) داریم که نمایانگر چهارگانه $(S, A, R, P)$ است \cite{RS_DRL_Article}. 

\subsubsection{تعریف فضای حالت}

فضای حالت مدل، وضعیت فعلی و پویای سیستم را با اندازه‌گیری سه معیار حیاتی عملکرد بازتاب می‌دهد. به بیان رسمی:
$$S = \{(E, P, L) \mid E \in \mathbb{R}, P \in [0,100], L \in [0,100]\}$$
که در آن:
\begin{itemize}
  \item \textbf{انرژی مصرفی ($E$):} مقدار حقیقیِ نمایانگر انرژی کل مصرفی توسط تمام سنسورهای شبکه در چرخه پیشین است.
  \item \textbf{تلفات بسته‌های شبکه ($P$):} درصدی از پیام‌های تولید شده که در طول مسیر به مقصد (دروازه اصلی) از دست رفته‌اند (بین ۰ تا ۱۰۰ درصد).
  \item \textbf{تاخیر ارسال ($L$):} معیاری معادل زمان مورد نیاز جهت رسیدن اطلاعات از سنسورها به گیت‌وی (حداکثر بازه ۱۰۰).
\end{itemize}

\subsubsection{فضای عمل: گزینه‌های وفقی گسسته}

فضای عمل شامل تمامی‌استراتژی‌ها و پیکربندی‌های ممکن به صورت مجموعه‌ای از گزینه‌های گسسته تعریف می‌شود:
$$A = \{a_1, a_2, \ldots, a_n\}$$
مقدار $n$ بسته به ابعاد شبکه \lr{DeltaIoT} برابر با ۲۱۶، ۱۲۹۶ یا ۴۰۹۶ در نظر گرفته می‌شود. هر عمل $a_i$ نمایانگر یک پیکربندی مشخص از پارامترهای $\{p_{i1}, p_{i2}, \ldots, p_{im}\}$ است.

\subsubsection{فرمول‌بندی تابع پاداش}

بر اساس ماهیت سیستم و اهداف، از فرمول‌بندی‌های زیر استفاده می‌شود:

\paragraph{پاداش چندهدفه:}
$$R = w_E \cdot \left(\frac{E_{max} - E}{E_{max} - E_{min}}\right) + w_P \cdot \left(\frac{P_{max} - P}{P_{max} - P_{min}}\right) + w_L \cdot \left(\frac{L_{max} - L}{L_{max} - L_{min}}\right)$$

\paragraph{پاداش تک‌هدفه:}
$$R = \left(\frac{E_{max} - E}{E_{max} - E_{min}}\right) \cdot w_E$$

\paragraph{اهداف آستانه‌ای (\lr{TTS} و \lr{TT}):}
تخصیص مقادیر پاداش مشروط به صورت دودویی:
$$R = \begin{cases} 
1 & \text{اگر } P < \theta_P \text{ و } L < \theta_L \\
0 & \text{در غیر این صورت}
\end{cases}$$

\section{بیان مسئله: یادگیری سیاست وفقی بهینه با عملیات مداوم}

پس از تعریف دقیق مؤلفه‌های \lr{MDP}، اکنون می‌توانیم مسئله اصلی بهینه‌سازی را که روش پیشنهادی \lr{RS-DRL} به دنبال حل آن است، به صورت صوری بیان کنیم. مسأله نهایی ما یادگیری سیاست وفقی بهینه $\pi^*: S \rightarrow A$ است که ارزش پاداش‌های تجمعی تنزیل‌یافته را حداکثر نماید:
$$\pi^* = \arg\max_{\pi} \mathbb{E}\left[\sum_{t=0}^{\infty} \gamma^t R(s_t, a_t) \mid \pi\right]$$

این سیاست باید همواره به‌طور مداوم در زمان اجرا عمل کرده و با چالش‌های عدم قطعیت محیطی، فضای عمل گسسته گسترده و پویایی غیرایستا (تغییر شرایط از یک چرخه به چرخه بعدی) مقابله نماید.

\section{خلاصه فصل}

این فصل چارچوب نظری و مدل‌سازی مورد نیاز برای درک رویکرد پیشنهادی را ترسیم نمود. ما با اصول سیستم‌های خودوفقی و الگوی \lr{MAPE-K} آغاز کردیم و سپس منابع عدم قطعیت و روش‌های مدل‌سازی زمان اجرا را بررسی نمودیم. با استفاده از \lr{DeltaIoT} به‌عنوان یک مطالعه موردی عینی، مسئله وفقی‌سازی را در قالب فرآیند تصمیم‌گیری مارکوف (\lr{MDP}) فرموله کردیم که شامل تعریف دقیق فضای حالت (انرژی، تلفات، تاخیر)، فضای عمل (تا ۴۰۹۶ گزینه) و توابع پاداش چندهدفه بود. در نهایت، هدف بهینه‌سازی را به‌عنوان یادگیری یک سیاست وفقی برای عملیات مداوم در زمان اجرا بیان کردیم. این زیربنای صوری و عملیاتی، امکان نقد و بررسی دقیق‌تر کارهای پیشین در فصل ۳ و ارائه متدولوژی نوآورانه \lr{RS-DRL} در فصل‌های آتی را فراهم می‌سازد.
\newpage
